\section{A finite harmonic reduction for every digamma power}
\label{hp:sec:all}

The newest source reports contain two independent fourth-power bridges.
The report \emph{Quartic Tails and Complex Resolvents}, Question~1 in its
research section, asks for the fifth power with every finite correction\cite{repo:quartic};
Question~2 asks about polynomials that remove lower-depth data.
\emph{Barnes Resolvents and Higher Harmonic Identities}, Question~F3, asks for
all higher powers and explicitly allows additional harmonic numerators\cite{repo:barnes}.
We answer the fifth-power problem and give a terminating reduction for
every fixed power. The latter has a precisely specified output: ordinary
convergent multiple zeta values and first Laurent coefficients of strict
multiple zeta functions. It does not assert that those coefficients
reduce to ordinary Stieltjes constants.

Write
\[
 H_N(s_1,\ldots,s_d)
   =\sum_{N\ge n_1>\cdots>n_d\ge1}
             n_1^{-s_1}\cdots n_d^{-s_d},\qquad
 Z_d(s_1,\ldots,s_d)=\lim_{N\to\infty}H_N(s_1,\ldots,s_d)
\]
when the latter converges. The larger index is first. The empty word
has value one, and $H_N^{(r)}=H_N(r)$, $H_N=H_N^{(1)}$.
The classical harmonic products, Euler-sum reductions and digamma
expansions keep their existing attribution\cite{dlmf:psi,hoffman,mezo}.
All indices in a harmonic word below are positive integers, except
for the one explicitly moving outer order. We use
\[
 \zeta(1+u)=u^{-1}+\sum_{j\ge0}\frac{(-1)^j\gamma_j}{j!}u^j,
 \qquad \gamma=\gamma_0,
 \qquad Q_k=\FP\int_0^1\psi(x)^k\,\dd x.
\]
Here $\FP$ is the constant in the cutoff coordinate $x=\varepsilon$:
it is the coefficient of $\varepsilon^0(\log\varepsilon)^0$ in
$\int_\varepsilon^1$. The upper endpoint is regular. Consequently
\[
 \FP\int_0^1x^{-j}\,\dd x=-\frac1{j-1}\quad(j\ge2),
 \qquad \FP\int_0^1x^{-1}\,\dd x=0.
\]

\subsection{The coefficients and the uniform theorem}

For $n\ge0$ introduce the formal series
\begin{equation}\label{hp:V}
 V_n(t)=-1+(H_n-\gamma)t
           +\sum_{r\ge2}\bigl(H_n^{(r)}+(-1)^r\zeta(r)\bigr)t^r.
\end{equation}
Only finitely many terms are needed in every operation that follows.
For $1\le j\le k$ put
\begin{equation}\label{hp:pdef}
 p_{k,j}(n)=[t^{k-j}]V_n(t)^k,
 \qquad \kappa_k=[t^k]V_0(t)^k,
 \qquad \Delta p_{k,j}(n)=p_{k,j}(n)-p_{k,j}(n-1).
\end{equation}
Define
\begin{align}
 B_k={}&\kappa_k+
   \sum_{j=2}^k\left\{
    \frac1{j-1}\sum_{n\ge1}\frac{\Delta p_{k,j}(n)}{n^{j-1}}
        -\sum_{n\ge1}\frac{p_{k,j}(n)}{n^j}\right\},
       \label{hp:Bdef}\\
 C_k={}&[L^0]\sum_{n=1}^N\frac{p_{k,1}(n)}n,
       \qquad L=\log N,\label{hp:Cdef}\\
 E_k(u)={}&\sum_{n\ge1}\Delta p_{k,1}(n)n^{-u}
       \qquad(\Re u>0).\label{hp:Edef}
\end{align}
The notation in \eqref{hp:Cdef} means the constant in the polynomial
asymptotic expansion in $L$, with an $o(1)$ remainder; existence is
part of the theorem. If $k=1$, the sum in \eqref{hp:Bdef} is empty.

\begin{theorem}[Uniform finite harmonic reduction]\label{hp:allpowers}
For every integer $k\ge1$, all series in \eqref{hp:Bdef} converge
absolutely, $C_k$ exists, $E_k$ is meromorphic near zero, and
\begin{equation}\label{hp:universal}
 \boxed{\displaystyle Q_k=B_k-C_k+[u^1]E_k(u).}
\end{equation}
Every term on the right has a finite algebraic construction from
\eqref{hp:V}. More precisely:
\begin{enumerate}
 \item $B_k$ and $C_k$ are finite linear combinations of convergent
 multiple zeta values, with coefficients in
 $\Q[\gamma,\zeta(2),\ldots,\zeta(k)]$. No limiting harmonic sum
 remains after the finite reduction given below.
 \item $E_k(u)$ is a finite linear combination of functions
 $Z_{d+1}(m+u,s_1,\ldots,s_d)$, where all displayed fixed indices are
 positive integers, $m+s_1+\cdots+s_d\le k-1$, and $d+1\le k-1$.
 Coefficients are in the same ordinary-constant algebra.
 \item For $k\ge2$, its only term of depth $k-1$ is
 \begin{equation}\label{hp:topdepth}
   -k!\,Z_{k-1}(1+u,\underbrace{1,\ldots,1}_{k-2}).
 \end{equation}
 All other moving-order terms have depth at most $k-2$.
\end{enumerate}
For $k=1$, $E_1=0$ and $Q_1=0$. For $k=2$, the formula gives
$Q_2=2\gamma_1-2\zeta(2)$.
\end{theorem}

The depth statement describes the specified finite normal form.
It is not an arithmetic independence statement. In particular,
shuffle, stuffle, duality, or an additional analytic identity may
further reduce the retained coordinates.

\subsection{The entire remainder and the normalized primitive}

\begin{lemma}[Mittag--Leffler identity at every power]\label{hp:ML}
Let $P_{k,n}(z)=\sum_{j=1}^kp_{k,j}(n)(n+z)^{-j}$. Away from the
nonpositive integers,
\begin{equation}\label{hp:MLformula}
 \psi(z)^k=P_{k,0}(z)+\kappa_k+
                    \sum_{n\ge1}\bigl(P_{k,n}(z)-P_{k,n}(0)\bigr).
\end{equation}
The series converges normally on compact subsets of that domain.
\end{lemma}

\begin{proof}
The digamma recurrence and its Taylor expansion at zero give
$t\psi(-n+t)=V_n(t)$. Therefore \eqref{hp:pdef} gives exactly the
principal parts, and $\kappa_k$ is the constant at the origin.
To determine the entire remainder, set $f(z)=\psi(z)^k-P_{k,0}(z)$
and integrate $zf(w)/(w(w-z))$ on $|w|=N+1/2$.
Uniformly on those circles, $f(w)=O(\log^kN)$.
In the right half-plane this is the usual digamma asymptotic.
In the left half-plane use
$\psi(w)=\psi(1-w)-\pi\cot(\pi w)$. The cotangent is uniformly
bounded on the chosen circles: where $|\Im w|\ge1/4$ its exponential
formula gives the bound; on the remaining arcs $\Re w$ stays a fixed
distance from the integers for sufficiently large $N$.
The contour integral is consequently $O(\log^kN/N)$ and tends to zero.
Its residues give \eqref{hp:MLformula}, including the constant
$f(0)=\kappa_k$.

Finally $p_{k,j}(n)=O(1+\log^{k-j}n)$ and
$(n+z)^{-j}-n^{-j}=O(n^{-j-1})$ on a fixed compact set.
This proves normal convergence. Matching principal parts alone
would not have proved the lemma: the vanishing contour integral
is what fixes the entire function.
\end{proof}

For $n\ge1$ set
\[
 F_1(n,z)=\Log(1+z/n)-z/n,\qquad
 F_j(n,z)=\frac{(n+z)^{1-j}-n^{1-j}}{1-j}-\frac{z}{n^j}
 \quad(j\ge2).
\]
On $\Re z>0$ use the logarithm real on the positive real axis.
The lemma immediately gives the following all-power antiderivative.

\begin{corollary}[A primitive with fixed additive constant]\label{hp:primitive}
For every $k\ge1$,
\begin{equation}\label{hp:primitiveformula}
 \mathcal P_k(z)=p_{k,1}(0)\Log z+
      \sum_{j=2}^k\frac{p_{k,j}(0)z^{1-j}}{1-j}
      +\kappa_kz+
      \sum_{n\ge1}\sum_{j=1}^kp_{k,j}(n)F_j(n,z)
\end{equation}
satisfies $\mathcal P_k'(z)=\psi(z)^k$ on $\Re z>0$.
Its expansion at the singular origin has constant zero. Thus
$\mathcal P_k(b)=\FP\int_0^b\psi(x)^k\,\dd x$ for $b>0$, and
ordinary integrals between positive endpoints are differences of
$\mathcal P_k$. Every fixed argument derivative of
\eqref{hp:MLformula} is also justified by normal convergence.
\end{corollary}

\begin{proof}
The displayed $F_j$ are the primitives, vanishing at zero, of
$(n+z)^{-j}-n^{-j}$. Their size on a fixed compact set is
$O(n^{-j-1})$; the same normal-convergence argument applies to
the differentiated series. On $|z|<1$ all nonsingular summands
vanish at zero. This proves the normalization as well as the
derivative identity.
\end{proof}

\subsection{Proof of the finite reduction}

Integrate \eqref{hp:MLformula}. For each $j\ge2$ finite summation by
parts and the vanishing terminal boundary term give
\begin{align*}
 &-\frac{p_{k,j}(0)}{j-1}
   +\sum_{n\ge1}p_{k,j}(n)
       \left\{\frac{n^{1-j}-(n+1)^{1-j}}{j-1}-n^{-j}\right\}\\
 &\hspace{15mm}=
   \frac1{j-1}\sum_{n\ge1}\frac{\Delta p_{k,j}(n)}{n^{j-1}}
         -\sum_{n\ge1}\frac{p_{k,j}(n)}{n^j}.
\end{align*}
All these sums are absolutely convergent. Hence
\begin{equation}\label{hp:afterintegration}
 Q_k=B_k+\sum_{n\ge1}p_{k,1}(n)
             \left\{\log(1+1/n)-1/n\right\}.
\end{equation}
Write $p(n)=p_{k,1}(n)$ and $d(n)=p(n)-p(n-1)$. At a finite cutoff,
\begin{equation}\label{hp:finiteparts}
 \sum_{n=1}^Np(n)\left\{\log(1+1/n)-1/n\right\}
 =p(N)\log(N+1)-\sum_{n=1}^Nd(n)\log n
                         -\sum_{n=1}^N\frac{p(n)}n.
\end{equation}
No divergent infinite rearrangement has been used.

The harmonic asymptotics imply, for a polynomial $P$,
$p(N)=P(\log N)+O(N^{-1}\log^M N)$. Thus the first term on the
right of \eqref{hp:finiteparts} has zero constant in $L=\log N$.
Likewise
\begin{equation}\label{hp:Dasym}
 d(n)=\frac{\sum_{r=0}^{k-2}a_r\log^r n}{n}
            +O(n^{-2}\log^M n)
\end{equation}
for some coefficients $a_r$ (the case of a zero polynomial is allowed).
Consequently
\[
 E_k(u)=\sum_ra_r(-1)^r\zeta^{(r)}(1+u)
   +\sum_{n\ge1}\left(d(n)-\frac{\sum_ra_r\log^r n}{n}\right)n^{-u},
\]
where the second sum is holomorphic for $\Re u>-1$, locally normally
with every fixed derivative. The Stieltjes coefficient convention gives
\begin{equation}\label{hp:coefficientconstant}
 [u^1]E_k(u)=-[L^0]\sum_{n=1}^Nd(n)\log n.
\end{equation}
Indeed $[u^1](-1)^r\zeta^{(r)}(1+u)=-\gamma_{r+1}$, while the
remaining coefficient is minus its absolutely convergent logarithmic
sum. Since \eqref{hp:afterintegration} converges, taking constants in
\eqref{hp:finiteparts} proves \eqref{hp:universal}.

Here are finite rules that compute every object in that proof. Expand
each harmonic polynomial by the strict quasi-shuffle rule
\begin{equation}\label{hp:stuffle}
 H_N(a)H_N(b,\mathbf v)
   =\sum_{\mathbf w\in(a)*(b,\mathbf v)}H_N(\mathbf w),
\end{equation}
where the recursive product inserts either initial letter or merges
them to their sum; repeated words retain their multiplicities. The
merge has a plus sign. For any nonempty word $\mathbf w=(w_1,\ldots,w_d)$,
\begin{align}
 \sum_{n=1}^N\frac{H_{n-1}(\mathbf w)}{n^j}
     &=H_N(j,\mathbf w),\label{hp:previous}\\
 \sum_{n=1}^N\frac{H_n(\mathbf w)}{n^j}
     &=H_N(j,\mathbf w)+H_N(j+w_1,w_2,\ldots,w_d).
       \label{hp:inclusive}
\end{align}
For the empty word the second term in \eqref{hp:inclusive} is absent.
The differences $\Delta p_{k,j}$ are computed by replacing
$H_n^{(r)}$ by $H_{n-1}^{(r)}+n^{-r}$ and subtracting.
Equations \eqref{hp:previous}--\eqref{hp:inclusive} then reduce $B_k$
to convergent multiple zeta values, because its initial index is
always at least two. They reduce $E_k$ to the finite moving-order
combination in the theorem.

For completeness, leading ones in $C_k$ can also be removed by an
explicit terminating rule. Let $\mathcal R_{\mathbf v}(T)$ denote
the harmonic regularization polynomial, with
\[
 \mathcal R_\varnothing(T)=1,\quad \mathcal R_{(1)}(T)=T,
 \quad\mathcal R_{\mathbf v}(T)=Z(\mathbf v)\quad(v_1\ge2).
\]
For $\mathbf v=(1,\mathbf w)$ with $a$ initial ones, write
$(1)*\mathbf w=a\mathbf v+\sum_{\mathbf u\ne\mathbf v}
m_{\mathbf u}\mathbf u$. Then
\begin{equation}\label{hp:regularization}
 \mathcal R_{\mathbf v}(T)=\frac1a
  \left\{T\mathcal R_{\mathbf w}(T)
       -\sum_{\mathbf u\ne\mathbf v}m_{\mathbf u}
                         \mathcal R_{\mathbf u}(T)\right\}.
\end{equation}
This decreases either total weight or the number of initial ones,
so it terminates. Before passing to $\mathcal R$, the same recurrence
with each terminal $Z(\mathbf v)$ replaced by $H_N(\mathbf v)$
is an exact finite identity, with $T=H_N$. Only afterward are the
convergent terminal words replaced by their limits. Their tails are
$O(N^{-1}\log^M N)$, and $H_N=L+\gamma+O(N^{-1})$.
Thus their constant in $L$ is obtained by $T=\gamma$.
This proves existence and the finite construction of $C_k$.
It also proves the individual continuation of every moving-order term:
when $m=1$, its inner harmonic word is a polynomial in $\log n$
plus $O(n^{-1}\log^M n)$, so a finite sum of zeta derivatives
subtracts its pole. When $m\ge2$, its defining series is already
holomorphic near $u=0$.

The index-weight bound follows directly from \eqref{hp:V}. Finally,
the unique harmonic monomial of degree $k-1$ in the simple-pole
coefficient is $-kH_n^{k-1}$. Its difference contributes
$-k(k-1)H_{n-1}^{k-2}/n$. The coefficient of the strict word of
$k-2$ ones in $H_{n-1}^{k-2}$ is $(k-2)!$.
Every other term has smaller harmonic degree after taking a
difference. This proves \eqref{hp:topdepth} and completes the theorem.

\section{The fifth-power identity and its additional spectral coordinate}
\label{hp:sec:fifth}

Define the inherited strict coefficients and the new height-one
coefficient by
\begin{equation}\label{hp:heightcoeffs}
 \eta=[u^1]Z_2(1+u,1),\qquad
 \theta=[u^1]Z_3(1+u,1,1),\qquad
 \phi=[u^1]Z_4(1+u,1,1,1).
\end{equation}
Brackets denote Laurent coefficients, with neither an additional
factorial nor a sign. A convergent depth-two derivative coordinate is
\begin{align}
 \chi_4&=2(\partial_1-\partial_2)Z_2(3,1)
                         +\partial_1Z_2(2,2)\label{hp:chi}\\
 &=-\sum_{n>m\ge1}\left\{
       \frac{2\log(n/m)}{n^3m}+\frac{\log n}{n^2m^2}\right\}.
       \notag
\end{align}
Both sums converge absolutely. The symbols $\phi$ and $\chi_4$ specify
the exact data retained below. There is no assertion that either is
arithmetically independent of the other constants.

\begin{theorem}[Fifth harmonic bridge]\label{hp:fifth}
The coordinate finite part of the fifth digamma power is
\begin{align}\label{hp:fifthmain}
 \boxed{\begin{aligned}
 Q_5={}&30\gamma^2\zeta(2)-60\gamma\zeta(3)+\frac{35}{3}\zeta(4)
       -120\phi+120\gamma\theta
       +60\bigl(\zeta(2)-\gamma^2\bigr)\eta\\
 &-20\bigl(\gamma^3-3\gamma\zeta(2)+2\zeta(3)\bigr)\gamma_1
       +10\chi_4-20\zeta'(4)
       +10\zeta(2)\zeta'(2)-20\gamma\zeta'(3).
 \end{aligned}}
\end{align}
An ordinary absolutely convergent series for the same constant is
\begin{equation}\label{hp:fifthseries}
 \boxed{\begin{aligned}
 Q_5={}&30\gamma^2\zeta(2)-60\gamma\zeta(3)+\frac{35}{3}\zeta(4)\\
      &+20\gamma_4-120\zeta(2)\gamma_2
       -5\sum_{n\ge1}\frac{\mathcal U(n)\log n}{n}.
 \end{aligned}}
\end{equation}
where, with $p=\psi(n)$ and $\ell=\log n$,
\begin{equation}\label{hp:Un}
 \mathcal U(n)=-4p^3+24\zeta(2)p-12p\psi'(n)-2\psi''(n)
    +4\ell^3-24\zeta(2)\ell
    +\frac{4\zeta(2)-2\psi'(n)}n+\frac{4p}{n^2}.
\end{equation}
The summand in \eqref{hp:fifthseries} is
$O(n^{-2}\log^3n)$.
\end{theorem}

\subsection{Explicit principal parts and every finite constant}

Put
\[
 b_n=H_n-\gamma,\quad c_n=H_n^{(2)}+\zeta(2),\quad
 d_n=H_n^{(3)}-\zeta(3),\quad e_n=H_n^{(4)}+\zeta(4).
\]
The principal coefficients supplied by \eqref{hp:pdef} are
\begin{equation}\label{hp:fifthparts}
 \begin{aligned}
 p_{5,5}&=-1,\qquad p_{5,4}=5b_n,\qquad p_{5,3}=5c_n-10b_n^2,\\
 p_{5,2}&=10b_n^3-20b_nc_n+5d_n,\qquad p_{5,1}=5r_n,\\
 r_n&=-b_n^4+6b_n^2c_n-2c_n^2-4b_nd_n+e_n.
 \end{aligned}
\end{equation}
The constant at zero is
\begin{equation}\label{hp:kappa5}
 \kappa_5=-\gamma^5+20\gamma^3\zeta(2)-30\gamma^2\zeta(3)
 -30\gamma\zeta(2)^2+20\gamma\zeta(4)
 +20\zeta(2)\zeta(3)-5\zeta(5).
\end{equation}
The exact difference of the residue is particularly simple:
\begin{equation}\label{hp:resdiff}
 r_n-r_{n-1}=\frac{4a_{n-1}}n+\frac{2c_{n-1}}{n^2}
                         +\frac{4b_{n-1}}{n^3},
 \qquad a_n=-b_n^3+3b_nc_n-d_n.
\end{equation}
It follows by substituting $b\mapsto b+n^{-1}$,
$c\mapsto c+n^{-2}$, $d\mapsto d+n^{-3}$, and
$e\mapsto e+n^{-4}$ into $r$. In particular, it is an exact finite
identity before any analytic continuation.

The lower-pole constant in \eqref{hp:Bdef} is
\begin{align}\label{hp:B5}
 B_5={}&-\gamma^5+30\gamma^3\zeta(2)+30\gamma^2\zeta(2)
      -80\gamma^2\zeta(3)-60\gamma\zeta(3)
      -\frac{65}{2}\gamma\zeta(4)\notag\\
 &\hspace{17mm}+60\zeta(2)\zeta(3)+\frac{35}{3}\zeta(4)-69\zeta(5).
\end{align}
We provide finite bookkeeping for this reduction and for the second
constant, so that the result is not dependent on a numerical fit.
The convergent values of weights at most five needed after the
quasi-shuffle expansion are
\begin{align}\label{hp:MZVtable}
 Z_2(2,1)&=\zeta(3),&Z_2(3,1)&=\tfrac14\zeta(4),&
 Z_2(2,2)&=\tfrac34\zeta(4),\notag\\
 Z_3(2,1,1)&=\zeta(4),&Z_4(2,1,1,1)&=\zeta(5),\notag\\
 Z_2(4,1)=Z_3(3,1,1)&=2\zeta(5)-\zeta(2)\zeta(3),\notag\\
 Z_2(3,2)=Z_3(2,2,1)&=3\zeta(2)\zeta(3)-\tfrac{11}{2}\zeta(5),\notag\\
 Z_2(2,3)=Z_3(2,1,2)&=\tfrac92\zeta(5)-2\zeta(2)\zeta(3).
\end{align}
These are classical Euler-sum and duality identities. For example,
Euler's formula for $Z_2(n,1)$ gives the $(4,1)$ value; the stuffle
product for $\zeta(2)\zeta(3)$ and its iterated-integral shuffle
product
\[
 \zeta(2)\zeta(3)=Z_2(2,3)+3Z_2(3,2)+6Z_2(4,1)
\]
give the two remaining depth-two values. Reversing the iterated
integral word gives the three indicated dualities. The height-one
values follow from the ordinary logarithmic beta integral. Also
$\zeta(2)^2=5\zeta(4)/2$.

For an additional check on \eqref{hp:B5}, let $B_{5,j}$ denote its
summand indexed by $j$ in \eqref{hp:Bdef}, excluding $\kappa_5$.
The finite rules above give
\begin{align}\label{hp:Bpieces}
 B_{5,2}={}&10\gamma^3\zeta(2)+30\gamma^2\zeta(2)
 -60\gamma^2\zeta(3)-70\gamma\zeta(3)+\tfrac{85}{2}\gamma\zeta(4)
 \notag\\
 &\quad+65\zeta(2)\zeta(3)+15\zeta(4)-\tfrac{215}{2}\zeta(5),\notag\\
 B_{5,3}={}&10\gamma^2\zeta(3)+10\gamma\zeta(3)-25\gamma\zeta(4)
       -30\zeta(2)\zeta(3)-5\zeta(4)+\tfrac{115}{2}\zeta(5),\notag\\
 B_{5,4}={}&5\gamma\zeta(4)+5\zeta(2)\zeta(3)
                            +\tfrac53\zeta(4)-15\zeta(5),\notag\\
 B_{5,5}={}&\zeta(5).
\end{align}
Adding these to \eqref{hp:kappa5} proves \eqref{hp:B5}.

To specify the constant of the simple-pole sum, set
\[
 A=4\gamma^3-12\gamma\zeta(2)+4\zeta(3),\quad
 C_0=-\gamma^4+6\gamma^2\zeta(2)-4\gamma\zeta(3)
                       -2\zeta(2)^2+\zeta(4),
 \quad e_j(N)=H_N(\underbrace{1,\ldots,1}_{j}).
\]
Expanding $r_n$ in strict words and applying
\eqref{hp:inclusive} gives the exact finite identity
\begin{align}\label{hp:finiteR}
 \sum_{n=1}^N\frac{r_n}{n}={}&-24e_5(N)+24\gamma e_4(N)
       +12(\zeta(2)-\gamma^2)e_3(N)+Ae_2(N)+C_0H_N\notag\\
 &-24H_N(2,1,1,1)+24\gamma H_N(2,1,1)
       +12(\zeta(2)-\gamma^2)H_N(2,1)+AH_N^{(2)}\notag\\
 &+4H_N(1,1,3)+2H_N(1,2,2)+4H_N(1,3,1)\notag\\
 &+4H_N(2,3)+2H_N(3,2)+4H_N(4,1)\notag\\
 &-4\gamma\{H_N(1,3)+H_N^{(4)}\}
          +2\zeta(2)\{H_N(1,2)+H_N^{(3)}\}.
\end{align}
One can apply \eqref{hp:regularization} directly. To make its
asymptotic evaluation especially short, write $T=H_N$ and note
\begin{align*}
 H_N(1,3)&=T\zeta(3)-\tfrac54\zeta(4)+o(1),\\
 H_N(1,2)&=T\zeta(2)-2\zeta(3)+o(1),\\
 4H_N(1,1,3)+2H_N(1,2,2)+4H_N(1,3,1)
   &=2T^2\zeta(3)-\tfrac52T\zeta(4)
                                  -2\zeta(2)\zeta(3)+o(1).
\end{align*}
Every $o(1)$ remains so after the displayed polynomial logarithmic
factors. The last line follows by stuffling $H_N$ into
$H_N(1,3)$, $H_N(2,2)$ and $H_N(3,1)$ and then using
\eqref{hp:MZVtable}. The elementary symmetric polynomials satisfy
\[
 e_j(N)=[t^j]\exp\left(Tt+
           \sum_{r=2}^j\frac{(-1)^{r-1}H_N^{(r)}}r t^r\right).
\]
Substitute these formulas and $T=L+\gamma+o(1)$ into
\eqref{hp:finiteR}. The complete result is
\begin{equation}\label{hp:Casym}
 \sum_{n=1}^N\frac{r_n}{n}
   =-\frac{L^5}{5}+4\zeta(2)L^3-18\zeta(4)L+C_r+o(1),
\end{equation}
where
\begin{equation}\label{hp:Cr}
 C_r=-\frac{\gamma^5}{5}+6\gamma^3\zeta(2)
      -16\gamma^2\zeta(3)-\frac{13}{2}\gamma\zeta(4)
      +12\zeta(2)\zeta(3)-\frac{69}{5}\zeta(5).
\end{equation}
Since $p_{5,1}=5r$, $C_5=5C_r$. Thus the substantial ordinary-constant
cancellation is exactly
\begin{equation}\label{hp:BCcancel}
 B_5-C_5=30\gamma^2\zeta(2)-60\gamma\zeta(3)+\frac{35}{3}\zeta(4).
\end{equation}
In particular, neither the finite constant at infinity nor the
finite parts of the higher poles have been discarded.

\subsection{The moving-order identity and the proof of the bridge}

Let $E_r(u)=E_5(u)/5$. From \eqref{hp:resdiff}, finite strict
harmonic identities give the full meromorphic germ
\begin{align}\label{hp:Er}
 E_r(u)={}&-24Z_4(1+u,1,1,1)+24\gamma Z_3(1+u,1,1)
         +12(\zeta(2)-\gamma^2)Z_2(1+u,1)\notag\\
 &+4(\gamma^3-3\gamma\zeta(2)+\zeta(3))\zeta(1+u)
       +4Z_2(1+u,3)\notag\\
 &+2Z_2(2+u,2)+2\zeta(2)\zeta(2+u)
                    +4Z_2(3+u,1)-4\gamma\zeta(3+u).
\end{align}
For example the identity used for its first line is
\[
 a_n=-6H_n(1,1,1)+6\gamma H_n(1,1)
       +3(\zeta(2)-\gamma^2)H_n
       +\gamma^3-3\gamma\zeta(2)+\zeta(3)+H_n^{(3)}.
\]
The series proof of \eqref{hp:Er} first applies in $\Re u>0$ and
then continues to the germ at zero. It is therefore legitimate to
extract its linear coefficient.

The only apparently singular additional depth-two term is removed
by an identity valid with a free outer order:
\[
 Z_2(1+u,3)=\zeta(1+u)\zeta(3)-\zeta(4+u)-Z_2(3,1+u).
\]
Consequently
\[
 [u^1]Z_2(1+u,3)
       =-\zeta(3)\gamma_1-\zeta'(4)-\partial_2Z_2(3,1).
\]
All remaining additional terms in \eqref{hp:Er} are ordinary
convergent derivatives. Extracting the coefficient and using
\eqref{hp:BCcancel} in \eqref{hp:universal} proves
\eqref{hp:fifthmain}.

For the independently useful convergent form, the elementary
polygamma relations give
\[
 b_{n-1}=\psi(n),\qquad c_{n-1}=2\zeta(2)-\psi'(n),\qquad
 d_{n-1}=\tfrac12\psi''(n).
\]
Substitution into \eqref{hp:resdiff} shows that
\[
 r_n-r_{n-1}
   =\frac{-4\log^3n+24\zeta(2)\log n+\mathcal U(n)}n.
\]
The standard positive-axis asymptotics of $\psi,\psi',\psi''$
give $\mathcal U(n)=O(n^{-1}\log^2 n)$.
The Stieltjes finite limits are
\[
 [L^0]\sum_{n=1}^N\frac{\log^j n}{n}=\gamma_j.
\]
Using them in \eqref{hp:coefficientconstant} proves
\eqref{hp:fifthseries}, with ordinary absolute convergence as
claimed.

\subsection{The highest-depth coefficient and a finite cancellation}

The normalization of $\phi$ has a direct elementary Mellin
description. In $\Re s>1$,
\begin{equation}\label{hp:heightMellin}
 Z_4(s,1,1,1)=\frac1{6\Gamma(s)}
   \int_0^\infty\frac{t^{s-1}[-\log(1-e^{-t})]^3}{e^t-1}\,\dd t.
\end{equation}
Subtract $(-\log t)^3/(6t)$ on $(0,1)$. The remainder is
holomorphic at $s=1$, and its value there is zero, because the
unsubtracted kernel is $-1/24$ times the derivative of
$[-\log(1-e^{-t})]^4$. Let $g_5=[u^5]\Gamma(1+u)^{-1}$.
Then
\begin{align}\label{hp:phiIntegral}
 \phi={}&g_5+\frac16\int_0^1\log t
  \left\{\frac{[-\log(1-e^{-t})]^3}{e^t-1}
                      -\frac{(-\log t)^3}{t}\right\}\,\dd t\notag\\
 &+\frac16\int_1^\infty
             \frac{\log t\,[-\log(1-e^{-t})]^3}{e^t-1}\,\dd t,
\end{align}
where both integrals converge ordinarily, and
\[
 g_5=\frac{\gamma^5-10\gamma^3\zeta(2)+20\gamma^2\zeta(3)
        +15\gamma\zeta(2)^2-30\gamma\zeta(4)
        -20\zeta(2)\zeta(3)+24\zeta(5)}{120}.
\]
This fixes the depth-four coordinate independently of a formal
Laurent label. It does not evaluate it in a smaller constant algebra.

The previously proved cubic and quartic identities\cite{repo:cubic,repo:quartic} are
\begin{align*}
 Q_3&=3\zeta(2)-6\eta-6\gamma\gamma_1,\\
 Q_4&=12\zeta(3)-12\gamma\zeta(2)+18\zeta(4)+4\zeta'(3)
       -24\theta+24\gamma\eta
       +12(\gamma^2-\zeta(2))\gamma_1.
\end{align*}
Combining them with the new theorem eliminates both inherited
strict coefficients:
\begin{corollary}[Eliminating the inherited height-one data]\label{hp:cancellation}
For $P_5(T)=T^5+5\gamma T^4+10(\gamma^2+\zeta(2))T^3$,
\begin{align}\label{hp:polyCancellation}
 \FP\int_0^1P_5(\psi(x))\,\dd x
 ={}&-120\phi+10\chi_4+
              \left(\frac{260}{3}+90\gamma\right)\zeta(4)
           -20\zeta'(4)+10\zeta(2)\zeta'(2)\notag\\
 &-20\bigl(\gamma^3+3\gamma\zeta(2)+2\zeta(3)\bigr)\gamma_1.
\end{align}
\end{corollary}
The extra derivative coordinate $\chi_4$ remains. Thus this is a
proved cancellation of $\eta$ and $\theta$, not a claimed solution
of the stronger all-power, single-coordinate polynomial problem.

\subsection{Verification and the exact scope of the result}

The accompanying code records sixty exact finite checks. It expands \eqref{hp:V} exactly, verifies
\eqref{hp:resdiff} and \eqref{hp:finiteR} on finite harmonic sums,
replays the terminating quasi-shuffle regularization, and reproduces
\eqref{hp:B5}--\eqref{hp:BCcancel}. Its numerical comparison uses
two analytically different methods: local integration of the Taylor
series of $(x\psi(x))^5$ followed by ordinary quadrature, and
\eqref{hp:fifthseries} with a finite asymptotic tail summed through
Hurwitz-zeta order derivatives. The direct evaluator gives
\[
 Q_5=-8.69030396262453826629558966236565357677029354492388429266\ldots.
\]
The recorded run agrees through forty-eight decimal places.
It is a floating-point diagnostic, not an interval certificate;
the exact proof is \eqref{hp:universal} together with the finite
coefficient calculations above. No numerical relation search is
used to infer a constant reduction.

The general theorem settles the existence of a finite reduction
for every power in its stated multiple-zeta-jet class and fixes
all entire and endpoint constants. The fifth-power theorem
supplies the explicit case requested in both newest reports.
The arithmetic reduction of $\chi_4$ or $\phi$, a proof that a
polynomial in $\psi$ can remove all lower-depth data at arbitrary
order, and a common-Hurwitz-shift counterpart remain separate
research problems.
